(1) 无边缘效应时介质覆盖面积为 \(wx\)，其余空气覆盖面积为 \(w(L-x)\)。这两个区域极板电势差相同，等效为并联电容，
\[
C(x)=\frac{\epsilon_0w(L-x)}d+\frac{\kappa\epsilon_0wx}d
=\frac{\epsilon_0w}d[L+(\kappa-1)x].
\]
所以 \(C'(x)=\epsilon_0w(\kappa-1)/d>0\)，介质插入必使总电容增加。

(2) 固定 \(Q\) 时与电源没有电量交换，全部电场静电储能为
\[
U_Q=\frac{Q^2}{2C(x)}.
\]
沿 \(x\) 方向由虚功关系求场对介质的力：
\[
F_Q=-\left(\frac{\partial U_Q}{\partial x}\right)_Q
=\frac{Q^2}{2C(x)^2}C'(x)>0.
\]
正方向是将介质继续吸入极板的方向，机械功由固定电荷体系的电场能减少提供。

(3) 在固定电压情况下，电荷 \(Q(x)=C(x)V\)，随介质插入而增加。电池在准静态插入 \(dx\) 时供给电功
\[
dW_{\rm bat}=V\,dQ=V^2\,dC,
\]
电容器的电场能增加 \(dU_{\rm field}=d(CV^2/2)=V^2dC/2\)。由能量守恒，另外的一半成为介质板的机械输出功：
\[
dW_{\rm mech}=dW_{\rm bat}-dU_{\rm field}
=\frac12V^2dC,
\quad
F_V=\frac12V^2C'(x)>0.
\]
亦可用包含电池的等效电势能 \(U_{\rm eff}=U_{\rm field}-VQ=-CV^2/2\) 计算。不能对固定 \(V\) 的 \(+CV^2/2\) 单独取负导数，因为那会错误地得到指向极板外的力。

(4) 从 \(x_1\) 到 \(x_2\)，固定 \(Q\) 时
\[
W_Q=\frac{Q^2}{2}\left[\frac1{C(x_1)}-\frac1{C(x_2)}\right]>0;
\]
固定 \(V\) 时
\[
W_V=\frac{V^2}{2}[C(x_2)-C(x_1)]
=\frac{\epsilon_0wV^2(\kappa-1)}{2d}(x_2-x_1)>0.
\]
当 \(\kappa\to1^+\) 时 \(C'(x)\to0\)，两种力与输出功都趋零，介质与空气不可区分。若要在有限实际极板中计算瞬态与端部力分布，需要补上本题忽略的边缘场。
